\documentclass[12pt, a4paper]{article}

% --- Pakiety ---
\usepackage[utf8]{inputenc}
\usepackage[T1]{fontenc}
\usepackage[polish]{babel}
\usepackage{amsmath, amssymb}
\usepackage{graphicx}
\usepackage{float}
\usepackage{geometry}
\usepackage{hyperref}
\usepackage{booktabs}
\usepackage{caption}
\usepackage{listings}
\usepackage{xcolor}
\usepackage{parskip}

\geometry{
  a4paper,
  left=2.5cm,
  right=2.5cm,
  top=2.5cm,
  bottom=2.5cm
}

% --- Zakaz dzielenia wyrazów ---
\hyphenpenalty=10000
\exhyphenpenalty=10000
\emergencystretch=3em

\hypersetup{
  colorlinks=true,
  linkcolor=blue,
  urlcolor=blue,
  citecolor=blue
}

% --- Zmiana nazwy float figure: Rysunek -> Wykres ---
\addto\captionspolish{\renewcommand{\figurename}{Wykres}}

% --- Styl listingów ---
\lstset{
  language=Python,
  basicstyle=\ttfamily\small,
  keywordstyle=\color{blue}\bfseries,
  commentstyle=\color{gray}\itshape,
  stringstyle=\color{orange!80!black},
  numbers=left,
  numberstyle=\tiny\color{gray},
  numbersep=8pt,
  frame=single,
  breaklines=true,
  breakatwhitespace=false,
  captionpos=b,
  showstringspaces=false,
  inputencoding=utf8,
  extendedchars=true,
  literate={ą}{{\k{a}}}1 {ę}{{\k{e}}}1 {ó}{{\'o}}1
            {ś}{{\'s}}1 {ł}{{\l{}}}1 {ź}{{\'z}}1
            {ż}{{\.z}}1 {ć}{{\'c}}1 {ń}{{\'n}}1
            {Ą}{{\k{A}}}1 {Ę}{{\k{E}}}1 {Ó}{{\'O}}1
            {Ś}{{\'S}}1 {Ł}{{\L{}}}1,
}

% --- Strona tytułowa ---
\title{
  \textbf{Metody Obliczeniowe w Nauce i Technice} \\[0.5em]
  \large Laboratorium 2 \textemdash{} Metoda najmniejszych kwadratów
}
\author{Mateusz Stoch}
\date{25 marca 2026}

\begin{document}

\maketitle

% ============================================================
\section{Zadanie 1 \textemdash{} Metoda najmniejszych kwadratów}
% ============================================================

% --- 1.1 Opis zadania ---
\subsection{Opis zadania}

Zadanie polega na zastosowaniu metody najmniejszych kwadratów do przewidywania
roku wydania utworu muzycznego na podstawie jego cech. Zbiór danych
\texttt{YearPredictionMSD.txt} zawiera $515\,345$ utworów, z~których pierwsze
$463\,715$ stanowi zbiór treningowy, a~pozostałe $51\,630$ to zbiór testowy.
Pierwsza kolumna to rok wydania, pozostałe $90$ kolumn to cechy numeryczne
(barwa i~kowariancja barwy).

Zadanie obejmuje:
\begin{enumerate}
  \item[(a)] Wczytanie zbioru i~podział na część treningową i~testową.
  \item[(b)] Narysowanie histogramu i~posortowanego wykresu roku wydania.
  \item[(c)--(d)] Zbudowanie macierzy modelu $\mathbf{A}$ (z dodaną kolumną jedynek)
        oraz jej znormalizowanej wersji metodą min-max do~$[0,1]$:
        \begin{equation}
          f_{\text{scaled}} = \frac{f - f_{\min}}{f_{\max} - f_{\min}}.
          \label{eq:minmax}
        \end{equation}
  \item[(e)] Wyznaczenie wektora wag $\mathbf{w}$ przez rozwiązanie równania normalnego:
        \begin{equation}
          \mathbf{A}^T \mathbf{A}\, \mathbf{w} = \mathbf{A}^T \mathbf{y}.
          \label{eq:normal}
        \end{equation}
  \item[(f)] Wyznaczenie wag metodą dekompozycji SVD (\texttt{scipy.linalg.lstsq})
        oraz z~regularyzacją Tichonowa ($\lambda = 0.01$):
        \begin{equation}
          \bigl(\mathbf{A}^T \mathbf{A} + \lambda\mathbf{I}\bigr)\mathbf{w}
          = \mathbf{A}^T \mathbf{y}.
          \label{eq:tikhonov}
        \end{equation}
  \item[(g)] Obliczenie współczynnika uwarunkowania macierzy $\mathbf{A}^T\mathbf{A}$
        dla wersji nieznormalizowanej i~znormalizowanej.
  \item[(h)] Ewaluacja modeli na zbiorze testowym przy użyciu średniego błędu
        bezwzględnego (MAE).
  \item[(i)] Inkrementalne wyznaczenie wag filtrem Kalmana podzielonym na
        5 pakietów danych.
\end{enumerate}

% --- 1.2 Argumentacja ---
\subsection{Argumentacja rozwiązania}

\subsubsection{Przygotowanie danych i normalizacja}

Do macierzy cech $\mathbf{X}$ dołączono z~lewej kolumnę jedynek, żeby model mógł
uczyć się wyrazu wolnego (biasu). Tak powstała macierz modelu $\mathbf{A}$.

Normalizacja min-max (wzór~\ref{eq:minmax}) jest potrzebna, bo poszczególne cechy mogą
mieć bardzo różne zakresy wartości -- bez normalizacji macierz
$\mathbf{A}^T\mathbf{A}$ staje się gorzej uwarunkowana i~rozwiązanie traci precyzję.
Analogicznie normalizuje się wektor $\mathbf{y}$. Żeby uzyskać predykcje w~oryginalnej
skali, stosuje się operację odwrotną:
\begin{equation}
  \hat{y}_{\text{orig}} = \hat{y}_{\text{scaled}}\cdot(y_{\max} - y_{\min}) + y_{\min}.
  \label{eq:unscale}
\end{equation}

\subsubsection{Równanie normalne i jego warianty}

Metoda najmniejszych kwadratów minimalizuje sumę kwadratów reszt
$\|\mathbf{A}\mathbf{w} - \mathbf{y}\|^2$. Rozwiązanie wyraża się przez równanie
normalne~(\ref{eq:normal}). Zamiast jawnie odwracać macierz, używamy solwera
\texttt{np.linalg.solve}, co jest numerycznie bezpieczniejsze.

Regularyzacja Tichonowa (wzór~\ref{eq:tikhonov}) dodaje do macierzy $\mathbf{A}^T\mathbf{A}$
człon $\lambda\mathbf{I}$, który poprawia uwarunkowanie i~chroni przed
przeuczeniem. Dla wyrazu wolnego (pierwsza kolumna jedynek) ustawiono
odpowiedni element macierzy $\mathbf{I}$ na zero, żeby nie regularyzować biasu.

Biblioteczna funkcja \texttt{scipy.linalg.lstsq} rozwiązuje problem wewnętrznie
przez rozkład SVD, co jest najbardziej stabilną, ale i~najwolniejszą metodą.

\subsubsection{Współczynnik uwarunkowania}

Współczynnik uwarunkowania $\operatorname{cond}(\mathbf{A}^T\mathbf{A})$ mówi,
o~ile maksymalnie mogą urosnąć błędy. Przybliżona utrata cyfr znaczących wynosi:
\begin{equation}
  \Delta \approx \log_{10}\!\bigl(\operatorname{cond}(\mathbf{A}^T\mathbf{A})\bigr).
\end{equation}

\subsubsection{Filtr Kalmana (wyznaczanie inkrementalne)}

Zamiast wczytać wszystkie dane naraz, można aktualizować wagi partiami.
Przy każdym pakiecie $\mathbf{A}_k$, $\mathbf{y}_k$ liczymy błąd predykcji
(innowację):
\begin{equation}
  \mathbf{e}_k = \mathbf{y}_k - \mathbf{A}_k\,\mathbf{w}_{k-1},
\end{equation}
aktualizujemy macierz informacji:
\begin{equation}
  \mathbf{P}_k^{-1} = \mathbf{P}_{k-1}^{-1} + \mathbf{A}_k^T\mathbf{A}_k,
\end{equation}
i poprawiamy wagi:
\begin{equation}
  \mathbf{w}_k = \mathbf{w}_{k-1}
    + \mathbf{P}_k\,\mathbf{A}_k^T\,\mathbf{e}_k.
\end{equation}
Wynik po pięciu iteracjach powinien zbiegnąć do tego, co daje rozwiązanie
jednorazowe. Zamiast jawnego odwracania $\mathbf{P}$ używamy \texttt{np.linalg.solve}.

\subsubsection{Kluczowe fragmenty kodu}

\begin{lstlisting}[caption={Budowa macierzy i normalizacja}]
y_train = train[0].values
X_train = train.drop(0, axis=1).values
y_test  = test[0].values
X_test  = test.drop(0, axis=1).values

ones_train = np.ones((X_train.shape[0], 1))
ones_test  = np.ones((X_test.shape[0],  1))
A_train = np.hstack((ones_train, X_train))
A_test  = np.hstack((ones_test,  X_test))

X_min = X_train.min(axis=0)
X_max = X_train.max(axis=0)
X_train_scaled = (X_train - X_min) / (X_max - X_min)
X_test_scaled  = (X_test  - X_min) / (X_max - X_min)

A_train_scaled = np.hstack((ones_train, X_train_scaled))
A_test_scaled  = np.hstack((ones_test,  X_test_scaled))

y_min = y_train.min();  y_max = y_train.max()
y_train_scaled = (y_train - y_min) / (y_max - y_min)
\end{lstlisting}

\begin{lstlisting}[caption={Wyznaczanie wag -- trzy warianty}]
# (e) Rownananie normalne -- solver
w_norm = np.linalg.solve(
    A_train.T @ A_train, A_train.T @ y_train)
w_norm_s = np.linalg.solve(
    A_train_scaled.T @ A_train_scaled,
    A_train_scaled.T @ y_train_scaled)

# (f-1) SVD przez lstsq
w_svd, _, _, _ = scipy.linalg.lstsq(A_train, y_train)

# (f-2) Regularyzacja Tichonowa lambda=0.01
lam = 0.01
I = np.eye(A_train.shape[1]);  I[0, 0] = 0
w_reg = np.linalg.solve(
    A_train.T @ A_train + lam * I, A_train.T @ y_train)
\end{lstlisting}

\begin{lstlisting}[caption={Filtr Kalmana -- aktualizacja partiami}]
P_inv    = 1e-5 * np.eye(A_train.shape[1])
w_kalman = np.zeros(A_train.shape[1])

batches = [(0,100000),(100000,200000),(200000,300000),
           (300000,400000),(400000,A_train.shape[0])]

for idx, (s, e) in enumerate(batches):
    A_k = A_train[s:e];  y_k = y_train[s:e]
    err   = y_k - A_k @ w_kalman
    P_inv += A_k.T @ A_k
    w_kalman += np.linalg.solve(P_inv, A_k.T @ err)
    y_pred = A_test @ w_kalman
    mae = np.mean(np.abs(y_test - y_pred))
    print(f"Iteracja {idx+1}: MAE = {mae:.4f}")
\end{lstlisting}

% --- 1.3 Wykresy i wyniki ---
\subsection{Wykresy i wyniki liczbowe}

\begin{figure}[H]
  \centering
  \includegraphics[width=\textwidth]{wykres1.png}
  \caption{Histogram (lewy) i~posortowany wykres roku wydania (prawy) dla całego zbioru.}
  \label{fig:wykres1}
\end{figure}

\begin{table}[H]
  \centering
  \caption{Współczynniki uwarunkowania macierzy $\mathbf{A}^T\mathbf{A}$.}
  \begin{tabular}{lc}
    \toprule
    Wersja macierzy & $\operatorname{cond}(\mathbf{A}^T\mathbf{A})$ \\
    \midrule
    Nieznormalizowana & $3{,}87 \times 10^{9}$ \\
    Znormalizowana    & $2{,}15 \times 10^{6}$ \\
    \bottomrule
  \end{tabular}
  \label{tab:cond}
\end{table}

\begin{table}[H]
  \centering
  \caption{Średni błąd bezwzględny (MAE) na zbiorze testowym dla różnych wariantów.}
  \begin{tabular}{lc}
    \toprule
    Metoda & MAE (lata) \\
    \midrule
    Równanie normalne (nieznormalizowane)  & $6{,}800$ \\
    Równanie normalne (znormalizowane)     & $6{,}800$ \\
    SVD -- \texttt{lstsq} (nieznormalizowane) & $6{,}800$ \\
    Regularyzacja Tichonowa $\lambda=0.01$ & $6{,}800$ \\
    \bottomrule
  \end{tabular}
  \label{tab:mae}
\end{table}

\begin{table}[H]
  \centering
  \caption{MAE filtra Kalmana po każdej iteracji (zbiór testowy).}
  \begin{tabular}{ccc}
    \toprule
    Iteracja & Wiersze treningowe & MAE (lata) \\
    \midrule
    1 & $0      - 100\,000$          & $6{,}846$ \\
    2 & $100\,000 - 200\,000$        & $6{,}807$ \\
    3 & $200\,000 - 300\,000$        & $6{,}800$ \\
    4 & $300\,000 - 400\,000$        & $6{,}801$ \\
    5 & $400\,000 - 463\,715$        & $6{,}800$ \\
    \bottomrule
  \end{tabular}
  \label{tab:kalman}
\end{table}

% --- 1.4 Wnioski ---
\subsection{Wnioski}

\begin{enumerate}
  \item Normalizacja min-max zmniejszyła współczynnik uwarunkowania
        z~$3{,}87 \times 10^9$ do $2{,}15 \times 10^6$, czyli o~ponad
        trzy rzędy wielkości (Tabela~\ref{tab:cond}). Przekłada się to bezpośrednio
        na mniejsze ryzyko utraty cyfr znaczących podczas rozwiązywania układu.

  \item Wszystkie cztery metody dały praktycznie identyczne MAE na poziomie
        $\approx 6{,}800$ lat (Tabela~\ref{tab:mae}). Oznacza to, że dla
        tego zbioru różnice numeryczne między solverem, SVD i~regularyzacją
        są na tyle małe, że nie wpływają na jakość predykcji.

  \item Filtr Kalmana już po 3 iteracjach (300\,000 wierszy) osiągnął wynik
        porównywalny z~rozwiązaniem jednorazowym~--- MAE $\approx 6{,}800$
        (Tabela~\ref{tab:kalman}). Podejście inkrementalne jest przydatne
        przy dużych zbiorach danych, które nie mieszczą się w~pamięci.

  \item Współczynnik uwarunkowania wpływa na liczbę traconych cyfr znaczących
        w~mantysie zgodnie z~zależnością $\Delta \approx \log_{10}(\text{cond})$.
        Dla nieznormalizowanej macierzy tracimy ok.~9 cyfr z~15 dostępnych
        w~podwójnej precyzji~--- dlatego normalizacja jest tu kluczowa.
\end{enumerate}

% ============================================================
\section{Podsumowanie}
% ============================================================

\begin{enumerate}
  \item \textbf{Normalizacja danych.}
        Przed rozwiązaniem układu warto znormalizować cechy, żeby macierz
        $\mathbf{A}^T\mathbf{A}$ była lepiej uwarunkowana.
        W~tym zadaniu uwarunkowanie poprawiło się o~ponad trzy rzędy wielkości.

  \item \textbf{Metody wyznaczania wag.}
        Równanie normalne z~solverem, SVD i~regularyzacja dają zbliżone wyniki
        dla dobrze przygotowanych danych. SVD jest najstabilniejszą metodą,
        ale też najwolniejszą. Regularyzacja przydaje się głównie przy danych,
        gdzie macierz może być prawie osobliwa.

  \item \textbf{Podejście inkrementalne.}
        Filtr Kalmana umożliwia aktualizację wag bez konieczności trzymania
        całego zbioru w~pamięci. Wynik zbiega do rozwiązania jednorazowego
        już po kilku iteracjach.
\end{enumerate}

% ============================================================
\section*{Bibliografia}
\addcontentsline{toc}{section}{Bibliografia}
% ============================================================

\begin{thebibliography}{9}

  \bibitem{lab02}
  \textit{Wprowadzenie do laboratorium}, materia\l{}y dost\k{e}pne na platformie
  Teams w~katalogu lab02.

\end{thebibliography}

\end{document}
